% ============================================================
% Sección 12: Herramientas de agentes
% ============================================================
\section{Herramientas de agentes}

\begin{frame}{¿Qué es una ``herramienta'' para un LLM?}
  \definicion{Tool / function calling}{
    Una función externa que el LLM puede invocar, descrita mediante un
    \textbf{schema} (nombre, descripción, parámetros) que el modelo usa
    para decidir cuándo y cómo llamarla.
  }
  \vfill
  \begin{itemize}
    \item El LLM \textbf{no ejecuta la función}: solo genera la
      \emph{intención} de llamarla (nombre + argumentos en JSON).
    \item El programa que envuelve al LLM es quien \textbf{ejecuta} la
      función real y le devuelve el resultado.
  \end{itemize}
\end{frame}

\begin{frame}[fragile]{Anatomía de un tool schema}
  \footnotesize
\begin{verbatim}
{
  "name": "calculadora",
  "description": "Evalua una expresion aritmetica",
  "parameters": {
    "type": "object",
    "properties": {
      "expresion": {"type": "string"}
    },
    "required": ["expresion"]
  }
}
\end{verbatim}
  \vfill
  \normalsize El modelo elige el nombre y llena los parámetros; tu código
  ejecuta \texttt{calculadora(expresion="3*7")}.
\end{frame}

\begin{frame}{Categorías de herramientas comunes}
  \small
  \begin{itemize}
    \item \textbf{Cómputo}: calculadora, intérprete de código (Python, SQL).
    \item \textbf{Información}: búsqueda web, RAG sobre documentos, APIs.
    \item \textbf{Sistema de archivos}: leer/escribir/editar archivos.
    \item \textbf{Acciones externas}: enviar correos, crear tickets, llamar APIs de terceros.
    \item \textbf{Memoria}: guardar/recuperar notas entre sesiones.
  \end{itemize}
\end{frame}

\begin{frame}{MCP: estandarizando las herramientas}
  \begin{itemize}
    \item \textbf{Model Context Protocol}: un servidor MCP expone un set de
      herramientas de forma estándar, consumible por cualquier cliente compatible.
    \item Ventaja: escribes la herramienta \textbf{una vez}, se conecta a
      Claude, a un IDE, o a cualquier agente que hable MCP.
    \item Es el equivalente a lo que USB hizo con los periféricos: un
      conector común en vez de uno distinto por dispositivo.
  \end{itemize}
\end{frame}
